% Generated draft: requires mathematical and source review.
\documentclass[11pt]{article}
\usepackage[margin=25mm]{geometry}
\usepackage{fontspec}
\setmainfont{Latin Modern Roman}
\usepackage{amsmath,amssymb,mathtools,unicode-math}
\setmathfont{Latin Modern Math}
\usepackage{enumitem,needspace,xurl,hyperref}
\setlength{\parindent}{0pt}
\setlength{\parskip}{5pt}
\setlength{\emergencystretch}{3em}
\newcommand{\N}{\mathbb{N}}
\newcommand{\Z}{\mathbb{Z}}
\newcommand{\Q}{\mathbb{Q}}
\newcommand{\R}{\mathbb{R}}
\newcommand{\C}{\mathbb{C}}
\newcommand{\F}{\mathbb{F}}
\newcommand{\A}{\mathbb{A}}
\newcommand{\PP}{\mathbb{P}}
\newcommand{\E}{\mathbb{E}}
\newcommand{\eps}{\varepsilon}
\newcommand{\dd}{\,\mathrm{d}}
\begin{document}
\section*{UnsolvedMath ID: 30001343}
\textbf{Source problem code:} catalogue entry
\subsection{Definitions and mathematical statement}
Let \((R,\mathfrak m)\) be a Noetherian excellent local domain, let \(k=R/\mathfrak m\) be algebraically closed, and assume that \(R\) is equicharacteristic, meaning that \(\operatorname{char}(R)=\operatorname{char}(k)\). Write \(X=\operatorname{Spec}R\) and let \(x\in X\) be its closed point.

The supplied question asks whether, for every such \(R\), there exist a finite-dimensional affine space \(\mathbb A_k^N\) and a toric birational map
\[
\pi:Z\longrightarrow\mathbb A_k^N
\]
such that the strict transform of \(X\) is nonsingular and transverse to every nondense toric orbit over \(x\). For a toric variety \(Z\) with dense torus \(T_Z\), a nondense toric orbit means an orbit \(O\ne T_Z\).

The source does not specify an embedding or other ambient realization of \(X\) in \(\mathbb A_k^N\). Consequently its strict transform under \(\pi\) is not defined by the supplied data. In a setting where an embedding \(X\hookrightarrow\mathbb A_k^N\) and a birational morphism \(\pi\) have been specified, and where \(U\subseteq\mathbb A_k^N\) is an open subset over which \(\pi\) is an isomorphism with \(X\cap U\) dense in \(X\), the strict transform would be the scheme-theoretic closure
\[
Y=\overline{\pi^{-1}(X\cap U)}^{\,Z}.
\]
This explains the terminology conditionally; it does not supply the missing embedding. The stated hypotheses do not assert that \(R\) is a finitely generated \(k\)-algebra, a localization of one, or equipped with a coefficient field. A finite-dimensional algebraic embedding must therefore not be assumed.

There are further unresolved qualifications. The source does not specify whether \(\pi\) must be a proper birational toric morphism, whether nonsingularity of \(Y\) is required everywhere or only above \(x\), or the precise meaning of transversality when the ambient toric variety is singular. In a smooth ambient setting, the usual tangent-space criterion for transversality at a point \(y\in Y\cap O\) is
\[
T_yY+T_yO=T_yZ,
\]
with \(Y\) and \(O\) smooth at \(y\). The phrase ``over the closed point'' suggests a condition at points mapping to \(x\), but the source does not state the exact quantifier. Thus the displayed question preserves the supplied target, while its ambient realization and these geometric conditions require clarification before it becomes a fully specified mathematical proposition.
\subsection{Short English statement}
For every excellent equicharacteristic Noetherian local domain with algebraically closed residue field, the source asks for a toric modification making its strict transform nonsingular and transverse to the toric boundary above the closed point. The ambient embedding and several geometric qualifications are unspecified, so this is a statement-only draft requiring mathematical review.
\subsection{Sources}
\begin{itemize}
\item Supplied local UnsolvedMath registry record, permanent ID \(30001343\), title ``Toric Embedded Resolution of Equicharacteristic Local Domains,'' source problem code \texttt{catalogue entry}, in the Oberwolfach Reports Open Problems collection. The record gives proposed year \(2009\) and names no proposer.
\item Supplied report citation: \emph{Singularities (2010)}, Oberwolfach Reports, pp.~2405--2470, DOI \texttt{10.4171/owr/2009/43}: \url{https://doi.org/10.4171/owr/2009/43}. This is also supplied as the literature reference ``original Oberwolfach report.'' The citation is reproduced as supplied; its bibliographic details and connection to the precise formulation have not been independently verified.
\item Supplied literature reference: ``Teissier on overweight deformations,'' \url{https://ems.press/books/ecr/197/3748}.
\item Supplied literature reference: ``one toric morphism for curves,'' \url{https://upcommons.upc.edu/bitstreams/bec04380-4ed3-45bf-8656-1ffdfda3c8dd/download}.
\item Literature assessment dated exactly \texttt{2026-08-21}: the supplied curation labels the problem ``partially solved'' and states that Teissier's construction is proved for important valuation and curve cases, while the stated arbitrary local-domain formulation remains open. This is supplied dated literature triage derived from the OWR curation, not an independently verified assessment; the status and cited claims require independent checking.
\item The supplied statement status is \texttt{corrected\_verified}. Its accompanying verification describes meaning-preserving semantic, typographic, and LaTeX cleanup without unsupported source-dependent additions. This registry label does not establish mathematical correctness or resolve the formulation issues above.
\end{itemize}
\end{document}
